% chapters/17/contoh.tex
\section{Dua Contoh Terapan}
\label{sec:contoh-bab-17}

\subsection{Contoh Pertama: Menghitung Nilai Akhir dari Komponennya}

\begin{contohbox}
Seorang mahasiswa memperoleh nilai sebagai berikut: tugas praktikum dan
portofolio 82, kuis 75, Ujian Tengah Semester 68, Mini Proyek AI 85, Proyek
Akhir 78, dan partisipasi 90. Berapa nilai akhirnya, dan bagaimana nilainya
tersebar pada keempat CPMK?
\end{contohbox}

Nilai akhir dihitung dari bobot komponen pada
Tabel~\ref{tab:komponen-nilai}:

\[
\begin{aligned}
\text{Nilai akhir} ={}& 0{,}25(82) + 0{,}10(75) + 0{,}20(68) \\
&+ 0{,}10(85) + 0{,}30(78) + 0{,}05(90) \\
={}& 20{,}5 + 7{,}5 + 13{,}6 + 8{,}5 + 23{,}4 + 4{,}5
= 78{,}0 .
\end{aligned}
\]

Nilai 78 termasuk huruf B pada Tabel~\ref{tab:konversi-nilai}.

Sebaran pada CPMK dihitung dari Tabel~\ref{tab:matriks-bobot}, dengan
membagi setiap sumbangan dengan bobot CPMK-nya. Untuk \CPMK{1}, yang berbobot
25 persen:

\[
\text{\CPMK{1}} = \frac{8(82) + 5(75) + 10(68) + 2(90)}{25}
= \frac{656 + 375 + 680 + 180}{25} = \frac{1\,891}{25} = 75{,}6 .
\]

Dengan cara yang sama, \CPMK{2} memperoleh
$(6 \times 82 + 5 \times 75 + 10 \times 68 + 10 \times 85 + 1 \times 90)/32
= 77{,}7$, \CPMK{3} memperoleh
$(6 \times 82 + 15 \times 78)/22 = 75{,}5$, dan \CPMK{4} memperoleh
$(5 \times 82 + 15 \times 78 + 1 \times 90)/21 = 79{,}5$.

Keempatnya di atas 70, sehingga seluruh CPMK tercapai. Ada satu hal yang tidak
tampak pada nilai akhir. \CPMK{1} dan \CPMK{3} hanya sedikit di atas ambang.
Keduanya menarik perhatian pada nilai UTS sebesar 68, yaitu nilai terendah, dan
UTS adalah satu-satunya instrumen yang mengukur \CPMK{1} dan \CPMK{2} tanpa
bantuan pekerjaan praktikum.

\subsection{Contoh Kedua: Menelusuri Kesenjangan Kompetensi}

\begin{contohbox}
Seorang mahasiswa memperoleh nilai praktikum 88, tetapi nilai UTS 55 dan nilai
ujian tulis pada UAS 58. Nilai akhirnya masih lulus, tetapi ia merasa tidak
memahami materinya. Ke mana ia harus menelusuri?
\end{contohbox}

Pola itu menunjukkan satu hal yang khas: pekerjaan praktikumnya baik,
sedangkan jawaban tulisnya lemah. Ada dua kemungkinan, dan keduanya perlu
diperiksa dengan cara yang berbeda.

\begin{langkah}
  \item \textbf{Periksa apakah praktikumnya dikerjakan sendiri.} Apabila
        pekerjaan praktikum dikerjakan bersama atau disalin, nilai tinggi pada
        praktikum tidak mencerminkan penguasaan, dan nilai tulis yang rendah
        adalah keterangan yang lebih jujur.
  \item \textbf{Periksa Sub-CPMK mana yang paling banyak ditanyakan pada
        UTS.} Pada Tabel~\ref{tab:peta-subcpmk}, UTS mengukur Sub-CPMK 1.1
        sampai 1.4. Apabila jawaban yang salah terpusat pada satu pokok
        bahasan, kesenjangannya ada di situ, bukan menyebar.
  \item \textbf{Periksa jenis pertanyaannya.} Soal UTS memuat dua bentuk:
        menjelaskan konsep dan memperbaiki kode. Nilai rendah pada bentuk
        pertama berarti pemahamannya belum sampai; nilai rendah pada bentuk
        kedua berarti ia dapat memakai konsepnya tetapi belum dapat
        menjelaskannya.
  \item \textbf{Kembali ke bab yang bersangkutan, dan kerjakan ulang
        latihannya.} Untuk Sub-CPMK 1.2, misalnya, kerjakan kembali latihan
        Bab~V dan jalankan kembali \code{ch05_gerak.py} dengan mengubah
        \istilah{delta time}-nya.
\end{langkah}

Tiga hal dapat dipetik dari contoh itu. Pertama, nilai akhir yang lulus tidak
berarti seluruh Sub-CPMK tercapai; ia hanya berarti jumlah berbobotnya cukup.
Kedua, sebaran nilai pada CPMK memperlihatkan \emph{di mana} kesenjangannya,
sedangkan nilai akhir menyembunyikannya. Ketiga, penelusuran selalu berakhir
pada satu bab tertentu, dan itulah gunanya peta pada
Tabel~\ref{tab:peta-subcpmk}.
